\chapter{Manual de usuario}

\section{Arranque}

Con la instalación ya preparada, el usuario levanta los servicios comunes,
el backend y el cliente web:

\begin{verbatim}
docker compose up -d
cd phase2-juan && uv run uvicorn simlab.api:app --port 8000
cd phase2-juan/frontend && npm run dev
\end{verbatim}

La interfaz queda disponible en \texttt{http://localhost:5173}. También
puede usarse la CLI con \texttt{uv run simlab}. La clave obligatoria es
\texttt{OPENROUTER\_API\_KEY}; las de Voyage AI y ZeroEntropy solo afectan
a la calidad de recuperación del Knowledge Backbone.

\section{Uso general}

El usuario describe en el chat el problema de toma de decisiones que quiere
estudiar. El Orchestrator presenta los modelos disponibles de la primera
fase, pide seleccionar uno o varios y delega la creación del entorno al
Architect. Tras aceptar la configuración, se ejecuta la simulación y la
interfaz muestra la rejilla, el \textit{replay}, los eventos críticos y el
estado de cada agente.

Al terminar, el Tracker resume las trayectorias, el Analyst interpreta los
patrones observados y el Reporter genera un informe PDF descargable. La
sesión queda registrada para consulta posterior.

En cualquier momento del proceso se puede iterar, pedir recomendaciones al
Orchestrator o volver a pasos anteriores. El Orchestrator se encarga de
guiar la conversación y de mantener la continuidad de la sesión.

\section{Interpretación de la interfaz}

La pantalla principal combina cuatro zonas. El chat permite pedir nuevas
simulaciones, aceptar propuestas del Orchestrator o solicitar explicaciones
sobre resultados ya generados. El panel de agentes muestra qué componente está
trabajando en cada momento: Architect para el entorno, Tracker para la
observación, Analyst para la interpretación y Reporter para el informe. La
rejilla de simulación muestra agentes y recursos, y permite reproducir la
ejecución paso a paso mediante el \textit{replay}. Las tarjetas de datos
resumen métricas, eventos críticos y resultados relevantes.

Las trazas de decisión deben leerse como una explicación local de cada paso.
Para cada agente, la interfaz muestra la percepción recibida, la acción
seleccionada, el resultado y, cuando el modelo los expone, sus valores internos
o alternativas valoradas. Estas trazas son útiles para distinguir dos casos:
un modelo que obtiene un mal resultado porque el entorno le fue desfavorable y
un modelo que obtiene un mal resultado porque su política de decisión eligió
acciones poco eficaces.

\section{Consulta de resultados}

Al terminar la ejecución, el informe PDF puede descargarse desde la propia
interfaz. El informe contiene la configuración del entorno, los modelos
comparados, las métricas principales, las gráficas generadas por el Analyst y
la interpretación textual de los patrones observados.

Las vistas del Knowledge Backbone permiten inspeccionar memorias, procedencia y
relaciones disponibles. En esta fase, las observaciones de simulación se
guardan en el \textit{namespace} \texttt{simulation}; los postulados y
formulaciones procedentes de la primera fase usan otros espacios de memoria.
Esta separación ayuda a distinguir entre conocimiento científico de entrada y
evidencia producida por ejecuciones del laboratorio.

\section{Errores frecuentes}

Si los puertos \texttt{8000} o \texttt{5173} están ocupados, se puede
levantar el backend en otro puerto y comunicarlo al frontend mediante
\texttt{VITE\_API\_PORT}. Si fallan PostgreSQL, MinIO, Neo4j o Qdrant, la
sesión no puede ejecutarse completa; normalmente basta con reiniciar
\texttt{docker compose up -d}. Si falta \texttt{OPENROUTER\_API\_KEY}, los
agentes no pueden invocarse.

Si no aparece ningún modelo disponible, conviene comprobar que la primera fase
ha registrado modelos aceptados en PostgreSQL y que los archivos asociados
existen en MinIO. Si el informe no se descarga, el problema suele estar en la
compilación LaTeX o en la ruta del artefacto persistido. Si el análisis parece
demasiado pobre, puede deberse a que Voyage AI o ZeroEntropy no estén
configurados y la recuperación del Knowledge Backbone se haya degradado.

\section{Términos básicos}

El Orchestrator es el agente coordinador de la sesión. Architect, Tracker,
Analyst y Reporter son agentes especializados. El Knowledge Backbone es la
capa compartida de memoria y persistencia. Un modelo de decisión es el
código Python generado por la primera fase que implementa \texttt{decide},
\texttt{update} y \texttt{get\_state}.
